\vspace{0.5em}\noindent\textbf{AWS COGNITO LÀ GÌ? TẠI SAO ĐÂY LÀ ``MẢNH GHÉP'' KHÔNG THỂ THIẾU
CHO HỆ THỐNG Y TẾ SỐ?}

Trong quá trình phát triển dự án Smart Healthcare Platform (Nền tảng hỗ
trợ sàng lọc y khoa AI và quản lý lịch hẹn), việc quản lý định danh và
phân quyền đa vai trò (Authentication \& Authorization) là một thách
thức lớn. Hệ thống phục vụ 3 nhóm người dùng hoàn toàn riêng biệt: Bệnh
nhân, Bác sĩ và Lễ tân. Thay vì tự phát triển một dịch vụ Auth thủ công
phức tạp, giải pháp tối ưu được lựa chọn là tích hợp AWS Cognito.

Các điểm chính cần nắm:

\begin{itemize}
\tightlist
\item
  AWS Cognito là dịch vụ quản lý định danh toàn diện (Fully Managed
  Identity \& Access Management), hỗ trợ cả User Pools (quản lý thư mục
  người dùng) và Identity Pools (phân quyền truy cập tài nguyên AWS).
\item
  Giải quyết triệt để bài toán phân quyền đa vai trò (RBAC - Role-Based
  Access Control) thông qua thuộc tính \texttt{cognito:groups} đính kèm
  trực tiếp trong JWT Tokens.
\item
  Hỗ trợ quy trình On-boarding an toàn cho nhân viên y tế: Admin khởi
  tạo tài khoản qua API, Cognito tự động gửi Email kèm mật khẩu tạm thời
  và bắt buộc đổi mật khẩu ở lần đăng nhập đầu tiên.
\item
  Backend Service (NestJS) xác thực Token nhẹ nhàng bằng Public Key của
  Cognito mà không cần lưu trữ bất kỳ mật khẩu nào của người dùng trong
  cơ sở dữ liệu.
\item
  Tích hợp chiến lược mã hóa UUID giúp ẩn danh định danh người dùng gốc,
  đáp ứng các tiêu chuẩn bảo mật dữ liệu y tế nhạy cảm.
\item
  Tối ưu chi phí phát triển và vận hành nhờ hạn mức AWS Free Tier cực
  rộng lên tới 50,000 MAU (Monthly Active Users) mỗi tháng.
\end{itemize}

Bài viết này giúp bạn hiểu rõ lý do tại sao nên nhường gánh nặng bảo mật
và hạ tầng Auth cho các dịch vụ Cloud-native như AWS Cognito để dồn
100\% nguồn lực vào việc tối ưu hóa logic nghiệp vụ cốt lõi của dự án.

Xem chi tiết
\href{https://www.facebook.com/groups/awsstudygroupfcj/permalink/2228021697962790/?rdid=MLFJXbxXfs8sMMKH\#}{Tại
đây}
